% ch7-theoreme-spectral.tex — Chapitre 7 : Théorème spectral et matrices symétriques
% Ce chapitre développe la théorie spectrale des opérateurs auto-adjoints :
% matrices symétriques et hermitiennes, diagonalisation orthogonale,
% opérateurs normaux, matrices définies positives, décomposition en valeurs
% singulières (SVD) et applications (ACP, formes quadratiques, optimisation).

\chapter{Théorème spectral et matrices symétriques}\label{ch:theoreme-spectral}

\section*{Introduction}
\addcontentsline{toc}{section}{Introduction}

Les matrices symétriques\index{matrice symétrique} occupent une place
privilégiée en algèbre linéaire~: elles apparaissent naturellement dans
de très nombreux contextes appliqués et jouissent de propriétés spectrales
remarquables qui les rendent particulièrement maniables.

\medskip

Voici quelques-unes de leurs occurrences les plus fréquentes~:
\begin{itemize}
  \item \textbf{Matrices de covariance~:}\index{matrice de covariance}
    en statistiques, la matrice de covariance
    $\Sigma = \mathbb{E}[(X-\mu)(X-\mu)^{\mathsf{T}}]$
    d'un vecteur aléatoire est symétrique (et semi-définie positive).
    L'analyse en composantes principales (ACP) repose entièrement sur
    sa diagonalisation.
  \item \textbf{Hessiennes~:}\index{hessienne} en optimisation, la matrice
    hessienne $H_f(x) = \bigl(\frac{\partial^2 f}{\partial x_i
    \partial x_j}\bigr)$ d'une fonction de classe $\mathcal{C}^2$
    est symétrique, et la nature de ses valeurs propres détermine la
    nature des points critiques.
  \item \textbf{Formes quadratiques~:}\index{forme quadratique} toute
    forme quadratique $q(\vec{x}) = \transpose{\vec{x}} A \vec{x}$
    peut être associée à une matrice symétrique~$A$, et sa classification
    (définie positive, semi-définie, indéfinie) est dictée par les valeurs
    propres de~$A$.
  \item \textbf{Vibrations et modes propres~:}\index{vibrations} en
    mécanique, l'équation $M\ddot{x} + Kx = 0$ conduit au problème
    aux valeurs propres généralisé $Kv = \omega^2 Mv$ où $K$ et $M$ sont
    symétriques.
  \item \textbf{Graphes et réseaux~:}\index{matrice laplacienne} la
    matrice laplacienne d'un graphe est symétrique et semi-définie positive~;
    ses valeurs propres gouvernent la connectivité et les processus de
    diffusion sur le graphe.
\end{itemize}

\medskip

Le résultat central de ce chapitre --- le \emph{théorème spectral}\index{théorème spectral} ---
affirme que toute matrice symétrique réelle (ou hermitienne complexe) est
\emph{orthogonalement diagonalisable}~: il existe une base orthonormée de
vecteurs propres. Ce théorème est l'un des résultats les plus importants
et les plus utilisés de toute l'algèbre linéaire.

\medskip

Dans tout ce chapitre, sauf mention contraire, $E$ désigne un espace
vectoriel de dimension finie $n \geq 1$ muni d'un produit scalaire
(réel) ou hermitien (complexe). On note $\inner{{\cdot},{\cdot}}$ ce
produit et $\norm{\cdot}$ la norme associée.

% ============================================================
\section{Opérateurs auto-adjoints}\label{sec:auto-adjoint}
% ============================================================

\begin{definition}{Adjoint d'un endomorphisme}{def:adjoint-ch7}
  Soit $E$ un espace préhilbertien\index{espace préhilbertien} de
  dimension finie et $f \in \End(E)$. L'\emph{adjoint}\index{adjoint}
  de $f$, noté $f^*$, est l'unique endomorphisme de $E$ vérifiant
  \[
    \forall \vec{u}, \vec{v} \in E, \qquad
    \inner{f(\vec{u}), \vec{v}} = \inner{\vec{u}, f^*(\vec{v})}.
  \]
\end{definition}

\begin{remark}{Existence et unicité de l'adjoint}{rem:adjoint-existence}
  L'existence de $f^*$ découle du théorème de représentation de Riesz
  appliqué à la forme linéaire $\vec{v} \mapsto \inner{f(\vec{u}), \vec{v}}$.
  L'unicité provient de la non-dégénérescence du produit scalaire~: si
  $\inner{\vec{u}, g_1(\vec{v})} = \inner{\vec{u}, g_2(\vec{v})}$ pour
  tout $\vec{u}$, alors $g_1 = g_2$.
\end{remark}

\begin{proposition}{Matrice de l'adjoint}{prop:mat-adjoint}
  Soit $\cB$ une base \emph{orthonormée} de~$E$ et
  $A = \Mat_\cB(f)$.
  \begin{enumerate}[label=(\roman*)]
    \item Sur $\R$~: $\Mat_\cB(f^*) = \transpose{A}$.
    \item Sur $\C$~: $\Mat_\cB(f^*) = \hermit{A} = \overline{\transpose{A}}$.
  \end{enumerate}
\end{proposition}

\begin{proof}
  Soient $(\vece{1},\ldots,\vece{n})$ la base orthonormée $\cB$, et
  $A = (a_{ij})$, $B = (b_{ij})$ les matrices de $f$ et $f^*$. On a
  \[
    b_{ij} = \inner{f^*(\vece{j}), \vece{i}}
           = \inner{\vece{j}, f(\vece{i})}
           = \overline{\inner{f(\vece{i}), \vece{j}}}
           = \overline{a_{ji}}.
  \]
  Donc $B = \hermit{A}$. Sur $\R$, $\overline{a_{ji}} = a_{ji}$, d'où
  $B = \transpose{A}$.
\end{proof}

\begin{definition}{Opérateur auto-adjoint (symétrique/hermitien)}{def:auto-adjoint}
  Un endomorphisme $f \in \End(E)$ est dit
  \emph{auto-adjoint}\index{opérateur auto-adjoint} (ou
  \emph{symétrique}\index{opérateur symétrique} si $\K = \R$,
  \emph{hermitien}\index{opérateur hermitien} si $\K = \C$) si
  $f^* = f$, c'est-à-dire
  \[
    \forall \vec{u}, \vec{v} \in E, \qquad
    \inner{f(\vec{u}), \vec{v}} = \inner{\vec{u}, f(\vec{v})}.
  \]
  En termes matriciels (dans une base orthonormée)~:
  \begin{itemize}
    \item sur $\R$~: $\transpose{A} = A$ (matrice symétrique),
    \item sur $\C$~: $\hermit{A} = A$ (matrice hermitienne).
  \end{itemize}
\end{definition}

\begin{example}{Matrices symétriques et hermitiennes}{ex:sym-herm}
  La matrice $A = \begin{pmatrix} 2 & -1 \\ -1 & 3 \end{pmatrix}$ est
  symétrique réelle~: $\transpose{A} = A$.

  La matrice $B = \begin{pmatrix} 3 & 2-i \\ 2+i & 5 \end{pmatrix}$ est
  hermitienne~: $\hermit{B} = B$. On note que les éléments diagonaux
  sont réels et que $b_{ji} = \overline{b_{ij}}$.
\end{example}

\begin{proposition}{Propriétés élémentaires des opérateurs auto-adjoints}{prop:auto-adj-props}
  Soit $f$ un opérateur auto-adjoint sur~$E$.
  \begin{enumerate}[label=(\roman*)]
    \item $\inner{f(\vec{v}), \vec{v}} \in \R$ pour tout $\vec{v} \in E$
      (même si $\K = \C$).
    \item Si $g$ est aussi auto-adjoint et $\alpha, \beta \in \R$, alors
      $\alpha f + \beta g$ est auto-adjoint.
    \item Si $f$ est inversible, alors $\inv{f}$ est auto-adjoint.
    \item $\Ker(f)$ et $\im(f)$ sont des sous-espaces orthogonaux~:
      $\Ker(f) = \im(f)^\perp$.
  \end{enumerate}
\end{proposition}

\begin{proof}\leavevmode
  \begin{enumerate}[label=(\roman*)]
    \item $\inner{f(\vec{v}), \vec{v}} = \inner{\vec{v}, f(\vec{v})}
      = \overline{\inner{f(\vec{v}), \vec{v}}}$, donc
      $\inner{f(\vec{v}), \vec{v}} \in \R$.
    \item $(\alpha f + \beta g)^* = \alpha f^* + \beta g^*
      = \alpha f + \beta g$ (car $\alpha, \beta \in \R$).
    \item $\inner{\inv{f}(\vec{u}), \vec{v}}
      = \inner{\inv{f}(\vec{u}), f(\inv{f}(\vec{v}))}
      = \inner{f(\inv{f}(\vec{u})), \inv{f}(\vec{v})}
      = \inner{\vec{u}, \inv{f}(\vec{v})}$.
    \item $\vec{v} \in \Ker(f) \iff f(\vec{v}) = \vzero
      \iff \inner{f(\vec{v}), \vec{u}} = 0 \;\forall \vec{u}
      \iff \inner{\vec{v}, f(\vec{u})} = 0 \;\forall \vec{u}
      \iff \vec{v} \in \im(f)^\perp$. \qedhere
  \end{enumerate}
\end{proof}

% ============================================================
\section{Valeurs propres réelles}\label{sec:vp-reelles}
% ============================================================

\begin{theorem}{Les valeurs propres d'un opérateur auto-adjoint sont réelles}{thm:vp-reelles}
  Soit $f$ un opérateur auto-adjoint sur un espace préhilbertien
  complexe\index{valeur propre!réelle} (ou euclidien) de dimension
  finie. Alors toutes les valeurs propres de $f$ sont réelles.
\end{theorem}

\begin{proof}
  Soit $\lambda \in \C$ une valeur propre de $f$ et $\vec{v} \neq \vzero$
  un vecteur propre associé~: $f(\vec{v}) = \lambda \vec{v}$.

  Calculons $\inner{f(\vec{v}), \vec{v}}$ de deux manières.

  D'une part~:
  \[
    \inner{f(\vec{v}), \vec{v}} = \inner{\lambda \vec{v}, \vec{v}}
    = \lambda \inner{\vec{v}, \vec{v}} = \lambda \norm{\vec{v}}^2.
  \]

  D'autre part, comme $f$ est auto-adjoint~:
  \[
    \inner{f(\vec{v}), \vec{v}} = \inner{\vec{v}, f(\vec{v})}
    = \inner{\vec{v}, \lambda \vec{v}}
    = \overline{\lambda} \inner{\vec{v}, \vec{v}}
    = \overline{\lambda} \norm{\vec{v}}^2.
  \]

  Puisque $\vec{v} \neq \vzero$, $\norm{\vec{v}}^2 > 0$, et en égalant~:
  \[
    \lambda \norm{\vec{v}}^2 = \overline{\lambda} \norm{\vec{v}}^2
    \implies \lambda = \overline{\lambda}
    \implies \lambda \in \R. \qedhere
  \]
\end{proof}

\begin{remark}{Nécessité du produit scalaire}{rem:necessite-ps}
  Ce résultat est faux sans hypothèse sur le produit scalaire~: la matrice
  $\begin{pmatrix} 0 & -1 \\ 1 & 0 \end{pmatrix}$ vérifie
  $\transpose{A} = -A$ (antisymétrique) et possède les valeurs propres
  $\pm i$. En revanche, pour une matrice symétrique réelle, la preuve
  ci-dessus s'applique en travaillant dans $\C^n$ muni du produit
  hermitien standard.
\end{remark}

% ============================================================
\section{Orthogonalité des vecteurs propres}\label{sec:vp-orthogonaux}
% ============================================================

\begin{theorem}{Orthogonalité des sous-espaces propres}{thm:vp-orthogonaux}
  Soit $f$ un opérateur auto-adjoint. Si $\lambda$ et $\mu$ sont deux
  valeurs propres \emph{distinctes}\index{vecteurs propres!orthogonaux}
  de $f$, alors les sous-espaces propres $E_\lambda(f)$ et $E_\mu(f)$
  sont orthogonaux~:
  \[
    \lambda \neq \mu \implies E_\lambda(f) \perp E_\mu(f).
  \]
\end{theorem}

\begin{proof}
  Soient $\vec{u} \in E_\lambda(f)$ et $\vec{v} \in E_\mu(f)$. Alors
  $f(\vec{u}) = \lambda \vec{u}$ et $f(\vec{v}) = \mu \vec{v}$.

  D'une part~:
  \[
    \inner{f(\vec{u}), \vec{v}} = \inner{\lambda \vec{u}, \vec{v}}
    = \lambda \inner{\vec{u}, \vec{v}}.
  \]
  (On utilise que $\lambda \in \R$ d'après le 
ef{thm:thm:vp-reelles}.)

  D'autre part, par auto-adjonction~:
  \[
    \inner{f(\vec{u}), \vec{v}} = \inner{\vec{u}, f(\vec{v})}
    = \inner{\vec{u}, \mu \vec{v}}
    = \overline{\mu} \inner{\vec{u}, \vec{v}}
    = \mu \inner{\vec{u}, \vec{v}}.
  \]

  En soustrayant~:
  \[
    (\lambda - \mu)\inner{\vec{u}, \vec{v}} = 0.
  \]
  Comme $\lambda \neq \mu$, on conclut $\inner{\vec{u}, \vec{v}} = 0$.
\end{proof}

% ============================================================
\section{Théorème spectral pour les matrices symétriques réelles}%
\label{sec:spectral-reel}
% ============================================================

\begin{theorem}{Théorème spectral réel}{thm:spectral-reel}
  Soit $A \in \Mat_n(\R)$ une matrice symétrique\index{théorème spectral!réel}
  ($\transpose{A} = A$). Alors~:
  \begin{enumerate}[label=(\roman*)]
    \item Toutes les valeurs propres de $A$ sont réelles.
    \item Il existe une matrice orthogonale\index{matrice orthogonale}
      $P \in \Orth(n)$ ($\transpose{P}P = I_n$) et une matrice diagonale
      $D = \diag(\lambda_1, \ldots, \lambda_n)$ telles que
      \[
        A = P D \transpose{P}.
      \]
    \item Les colonnes de $P$ forment une base orthonormée de $\R^n$
      constituée de vecteurs propres de~$A$.
  \end{enumerate}
  De manière intrinsèque~: tout opérateur auto-adjoint sur un espace
  euclidien de dimension finie admet une base orthonormée de vecteurs
  propres.
\end{theorem}

\begin{proof}
  La preuve procède par récurrence\index{récurrence} forte sur la
  dimension $n$.

  \medskip
  \noindent\emph{Initialisation~:} pour $n = 1$, toute matrice $A = (a)$
  est triviallement symétrique et diagonale~; on prend $P = (1)$ et
  $D = (a)$.

  \medskip
  \noindent\emph{Hérédité~:} supposons le résultat vrai pour tout espace
  euclidien de dimension $< n$, et soit $f$ un opérateur auto-adjoint sur
  un espace euclidien $E$ de dimension~$n$.

  \medskip
  \noindent\textbf{Étape 1~: Existence d'une valeur propre réelle.}

  Considérons la matrice $A = \Mat_\cB(f)$ dans une base orthonormée
  $\cB$. Comme $\transpose{A} = A$, on peut voir $A$ comme une matrice
  hermitienne dans $\C^n$. Le polynôme caractéristique
  $\chi_A(\lambda) = \det(A - \lambda I_n)$ est de degré $n$ à
  coefficients réels. Par le théorème fondamental de l'algèbre, il
  possède au moins une racine $\lambda_1 \in \C$. Par le
  
ef{thm:thm:vp-reelles}, $\lambda_1 \in \R$.

  \medskip
  \noindent\textbf{Étape 2~: Choix d'un vecteur propre unitaire.}

  Soit $\vec{v}_1 \in E$ un vecteur propre associé à $\lambda_1$ avec
  $\norm{\vec{v}_1} = 1$.

  \medskip
  \noindent\textbf{Étape 3~: Invariance du complément orthogonal.}

  Posons $W = (\Vect(\vec{v}_1))^\perp = \set{\vec{w} \in E :
  \inner{\vec{w}, \vec{v}_1} = 0}$. Montrons que $W$ est stable par $f$.

  Soit $\vec{w} \in W$. Alors~:
  \[
    \inner{f(\vec{w}), \vec{v}_1}
    = \inner{\vec{w}, f(\vec{v}_1)}
    = \inner{\vec{w}, \lambda_1 \vec{v}_1}
    = \lambda_1 \inner{\vec{w}, \vec{v}_1} = 0.
  \]
  Donc $f(\vec{w}) \in W$~: le sous-espace $W$ est invariant par~$f$.

  \medskip
  \noindent\textbf{Étape 4~: Application de l'hypothèse de récurrence.}

  La restriction $g = f\restrict{W} \in \End(W)$ est auto-adjointe sur
  l'espace euclidien $W$ (le produit scalaire est hérité de~$E$) de
  dimension $n - 1$. Par hypothèse de récurrence, il existe une base
  orthonormée $(\vec{v}_2, \ldots, \vec{v}_n)$ de $W$ constituée de
  vecteurs propres de~$g$ (donc de~$f$), associés à des valeurs propres
  réelles $\lambda_2, \ldots, \lambda_n$.

  \medskip
  \noindent\textbf{Étape 5~: Conclusion.}

  La famille $(\vec{v}_1, \vec{v}_2, \ldots, \vec{v}_n)$ est une base
  orthonormée de $E$ (car $\vec{v}_1 \perp W$ et
  $(\vec{v}_2, \ldots, \vec{v}_n)$ est orthonormée dans~$W$). Chaque
  $\vec{v}_i$ est vecteur propre de $f$ pour $\lambda_i$.

  En posant $P = (\vec{v}_1 \mid \vec{v}_2 \mid \cdots \mid \vec{v}_n)$,
  la matrice dont les colonnes sont les $\vec{v}_i$, on obtient
  $P \in \Orth(n)$ et $\transpose{P} A P = D = \diag(\lambda_1, \ldots,
  \lambda_n)$, soit $A = P D \transpose{P}$.
\end{proof}

\begin{corollary}{Existence d'une base orthonormée de vecteurs propres}{cor:bon-vp}
  Soit $A \in \Sym_n(\R)$. Alors $\R^n$ possède une base orthonormée
  formée de vecteurs propres de~$A$, et le polynôme caractéristique
  de $A$ est scindé sur~$\R$.
\end{corollary}

\begin{remark}{Interprétation géométrique}{rem:spectral-geometrique}
  L'écriture $A = P D \transpose{P}$ signifie que l'application linéaire
  associée à $A$ agit par~:
  \begin{enumerate}[label=(\arabic*)]
    \item rotation (passage en coordonnées dans la base propre via
      $\transpose{P}$),
    \item dilatation le long des axes propres (multiplication par les
      valeurs propres via~$D$),
    \item rotation inverse (retour à la base de départ via~$P$).
  \end{enumerate}
\end{remark}

% ============================================================
\section{Théorème spectral pour les matrices hermitiennes}%
\label{sec:spectral-hermitien}
% ============================================================

\begin{theorem}{Théorème spectral hermitien}{thm:spectral-hermitien}
  Soit $A \in \Mat_n(\C)$ une matrice hermitienne\index{théorème spectral!hermitien}
  ($\hermit{A} = A$). Alors~:
  \begin{enumerate}[label=(\roman*)]
    \item Toutes les valeurs propres de $A$ sont réelles.
    \item Il existe une matrice unitaire\index{matrice unitaire}
      $U \in \Mat_n(\C)$ ($\hermit{U}U = I_n$) et une matrice diagonale
      réelle $D = \diag(\lambda_1, \ldots, \lambda_n)$ telles que
      \[
        A = U D \hermit{U}.
      \]
    \item Les colonnes de $U$ forment une base orthonormée de $\C^n$
      (pour le produit hermitien standard) constituée de vecteurs propres
      de~$A$.
  \end{enumerate}
\end{theorem}

\begin{proof}
  La preuve est identique à celle du 
ef{thm:thm:spectral-reel}, en
  remplaçant «~orthogonal~» par «~unitaire~», «~transposée~» par
  «~transconjuguée~» et le produit scalaire euclidien par le produit
  hermitien. L'existence d'une valeur propre est immédiate car $\C$ est
  algébriquement clos, et la réalité des valeurs propres est assurée par
  le 
ef{thm:thm:vp-reelles}. L'argument d'invariance du complément
  orthogonal reste valide~: si $f(\vec{v}_1) = \lambda_1 \vec{v}_1$
  avec $\lambda_1 \in \R$ et $\vec{w} \perp \vec{v}_1$, alors
  $\inner{f(\vec{w}), \vec{v}_1} = \inner{\vec{w}, f(\vec{v}_1)}
  = \overline{\lambda_1}\inner{\vec{w}, \vec{v}_1} = 0$.
\end{proof}

% ============================================================
\section{Diagonalisation orthogonale~: procédure et exemples}%
\label{sec:diag-ortho}
% ============================================================

La procédure de \emph{diagonalisation orthogonale}\index{diagonalisation orthogonale}
d'une matrice symétrique réelle $A \in \Sym_n(\R)$ se déroule en
quatre étapes~:

\begin{enumerate}[label=\textbf{Étape \arabic*}.]
  \item \textbf{Polynôme caractéristique~:} calculer
    $\chi_A(\lambda) = \det(A - \lambda I_n)$ et trouver ses racines
    $\lambda_1, \ldots, \lambda_r$ (avec multiplicités).
  \item \textbf{Sous-espaces propres~:} pour chaque $\lambda_i$, déterminer
    $E_{\lambda_i} = \Ker(A - \lambda_i I_n)$ par résolution du système
    homogène.
  \item \textbf{Orthonormalisation~:} si $\dim E_{\lambda_i} > 1$,
    appliquer le procédé de Gram--Schmidt\index{Gram--Schmidt} pour obtenir
    une base orthonormée de chaque sous-espace propre. (Les vecteurs
    propres associés à des valeurs propres \emph{distinctes} sont déjà
    orthogonaux par le 
ef{thm:thm:vp-orthogonaux}.)
  \item \textbf{Assemblage~:} former $P$ en juxtaposant les vecteurs
    propres orthonormés, et $D = \transpose{P} A P$.
\end{enumerate}

\begin{example}{Diagonalisation orthogonale d'une matrice $2 \times 2$}{ex:diag-ortho-2x2}
  Soit $A = \begin{pmatrix} 5 & 2 \\ 2 & 2 \end{pmatrix}$.

  \textbf{Étape 1.} $\chi_A(\lambda) = (5-\lambda)(2-\lambda) - 4
  = \lambda^2 - 7\lambda + 6 = (\lambda - 1)(\lambda - 6)$.
  Valeurs propres~: $\lambda_1 = 1$, $\lambda_2 = 6$.

  \textbf{Étape 2.}
  \begin{itemize}
    \item $E_1 = \Ker(A - I_2) = \Ker\begin{pmatrix} 4 & 2 \\ 2 & 1
      \end{pmatrix} = \Vect\begin{pmatrix} 1 \\ -2 \end{pmatrix}$.
    \item $E_6 = \Ker(A - 6I_2) = \Ker\begin{pmatrix} -1 & 2 \\ 2 & -4
      \end{pmatrix} = \Vect\begin{pmatrix} 2 \\ 1 \end{pmatrix}$.
  \end{itemize}

  \textbf{Étape 3.} On normalise~:
  $\vec{v}_1 = \frac{1}{\sqrt{5}}\begin{pmatrix} 1 \\ -2 \end{pmatrix}$,
  $\vec{v}_2 = \frac{1}{\sqrt{5}}\begin{pmatrix} 2 \\ 1 \end{pmatrix}$.
  On vérifie~: $\inner{\vec{v}_1, \vec{v}_2} = \frac{1}{5}(2-2) = 0$.

  \textbf{Étape 4.}
  $P = \frac{1}{\sqrt{5}}\begin{pmatrix} 1 & 2 \\ -2 & 1 \end{pmatrix}$,
  $D = \begin{pmatrix} 1 & 0 \\ 0 & 6 \end{pmatrix}$, et
  $A = P D \transpose{P}$.
\end{example}

\begin{example}{Diagonalisation orthogonale d'une matrice $3 \times 3$}{ex:diag-ortho-3x3}
  Soit $A = \begin{pmatrix} 2 & 1 & 1 \\ 1 & 2 & 1 \\ 1 & 1 & 2 \end{pmatrix}$.

  \textbf{Étape 1.} On calcule~:
  \[
    \chi_A(\lambda) = -(\lambda - 4)(\lambda - 1)^2.
  \]
  Valeurs propres~: $\lambda_1 = 4$ (multiplicité 1),
  $\lambda_2 = 1$ (multiplicité 2).

  \textbf{Étape 2.}
  \begin{itemize}
    \item $E_4 = \Ker(A - 4I_3)
      = \Vect\begin{pmatrix} 1 \\ 1 \\ 1 \end{pmatrix}$.
    \item $E_1 = \Ker(A - I_3)
      = \setcond{\begin{pmatrix} x \\ y \\ z \end{pmatrix}}{x + y + z = 0}
      = \Vect\left(\begin{pmatrix} -1 \\ 1 \\ 0 \end{pmatrix},\;
      \begin{pmatrix} -1 \\ 0 \\ 1 \end{pmatrix}\right)$.
  \end{itemize}

  \textbf{Étape 3.} Pour $E_4$~: $\vec{v}_1
  = \frac{1}{\sqrt{3}}\begin{pmatrix} 1 \\ 1 \\ 1 \end{pmatrix}$.

  Pour $E_1$, on applique Gram--Schmidt.
  Posons $\vec{w}_1 = \begin{pmatrix} -1 \\ 1 \\ 0 \end{pmatrix}$,
  $\vec{w}_2 = \begin{pmatrix} -1 \\ 0 \\ 1 \end{pmatrix}$.

  $\vec{u}_1 = \vec{w}_1$, $\norm{\vec{u}_1} = \sqrt{2}$.

  $\vec{u}_2 = \vec{w}_2 - \frac{\inner{\vec{w}_2, \vec{u}_1}}{\norm{\vec{u}_1}^2}\vec{u}_1
  = \begin{pmatrix} -1 \\ 0 \\ 1 \end{pmatrix}
  - \frac{1}{2}\begin{pmatrix} -1 \\ 1 \\ 0 \end{pmatrix}
  = \begin{pmatrix} -1/2 \\ -1/2 \\ 1 \end{pmatrix}$.
  On a $\norm{\vec{u}_2} = \sqrt{3/2}$.

  Après normalisation~:
  $\vec{v}_2 = \frac{1}{\sqrt{2}}\begin{pmatrix} -1 \\ 1 \\ 0 \end{pmatrix}$,
  $\vec{v}_3 = \frac{1}{\sqrt{6}}\begin{pmatrix} -1 \\ -1 \\ 2 \end{pmatrix}$.

  \textbf{Étape 4.}
  \[
    P = \begin{pmatrix}
      1/\sqrt{3} & -1/\sqrt{2} & -1/\sqrt{6} \\
      1/\sqrt{3} & 1/\sqrt{2}  & -1/\sqrt{6} \\
      1/\sqrt{3} & 0           & 2/\sqrt{6}
    \end{pmatrix}, \qquad
    D = \begin{pmatrix} 4 & 0 & 0 \\ 0 & 1 & 0 \\ 0 & 0 & 1 \end{pmatrix}.
  \]
\end{example}

% ============================================================
\section{Opérateurs normaux}\label{sec:normaux}
% ============================================================

\begin{definition}{Opérateur normal}{def:normal}
  Un endomorphisme $f \in \End(E)$ est dit
  \emph{normal}\index{opérateur normal} si $f$ commute avec son adjoint~:
  \[
    f \circ f^* = f^* \circ f.
  \]
  En termes matriciels (dans une base orthonormée), $A$ est
  normale\index{matrice normale} si $\hermit{A}A = A\hermit{A}$.
\end{definition}

\begin{remark}{Exemples d'opérateurs normaux}{rem:ex-normaux}
  Les classes suivantes sont toutes constituées d'opérateurs normaux~:
  \begin{itemize}
    \item opérateurs auto-adjoints ($f^* = f$),
    \item opérateurs anti-auto-adjoints ($f^* = -f$),
    \item opérateurs unitaires/orthogonaux ($f^* = \inv{f}$).
  \end{itemize}
\end{remark}

\begin{lemma}{Caractérisation par la norme}{lem:normal-norme}
  $f$ est normal si et seulement si
  $\norm{f(\vec{v})} = \norm{f^*(\vec{v})}$ pour tout $\vec{v} \in E$.
\end{lemma}

\begin{proof}
  On a $\norm{f(\vec{v})}^2 = \inner{f(\vec{v}), f(\vec{v})}
  = \inner{f^*f(\vec{v}), \vec{v}}$ et de même
  $\norm{f^*(\vec{v})}^2 = \inner{ff^*(\vec{v}), \vec{v}}$.
  Donc $\norm{f(\vec{v})}^2 - \norm{f^*(\vec{v})}^2
  = \inner{(f^*f - ff^*)(\vec{v}), \vec{v}}$.

  Si $f$ est normal, $f^*f = ff^*$ et l'expression est nulle.
  Réciproquement, si $\inner{g(\vec{v}), \vec{v}} = 0$ pour tout
  $\vec{v}$ avec $g = f^*f - ff^*$ auto-adjoint, alors $g = 0$ (par
  polarisation), donc $f^*f = ff^*$.
\end{proof}

\begin{theorem}{Théorème spectral pour les opérateurs normaux sur $\C$}{thm:spectral-normal}
  Soit $E$ un espace hermitien de dimension finie et
  $f \in \End(E)$.\index{théorème spectral!opérateurs normaux}
  Les assertions suivantes sont équivalentes~:
  \begin{enumerate}[label=(\roman*)]
    \item $f$ est normal~: $f \circ f^* = f^* \circ f$.
    \item Il existe une base orthonormée de $E$ constituée de vecteurs
      propres de~$f$.
  \end{enumerate}
  En termes matriciels~: $A \in \Mat_n(\C)$ est normale si et seulement
  s'il existe une matrice unitaire $U$ telle que
  $A = U \diag(\lambda_1, \ldots, \lambda_n) \hermit{U}$.
\end{theorem}

\begin{proof}\leavevmode

  \noindent$(\mathrm{ii}) \Rightarrow (\mathrm{i})$~:
  Si $A = U D \hermit{U}$ avec $D$ diagonale, alors
  $\hermit{A} = U \overline{D} \hermit{U}$, et
  $A\hermit{A} = U D\overline{D} \hermit{U}
  = U \overline{D}D \hermit{U} = \hermit{A}A$ car $D\overline{D}
  = \overline{D}D$ (matrices diagonales commutent).

  \medskip
  \noindent$(\mathrm{i}) \Rightarrow (\mathrm{ii})$~:
  Par récurrence sur $n = \dim E$.

  Pour $n = 1$, c'est trivial. Supposons le résultat pour
  $n - 1$ et soit $f$ normal sur $E$ de dimension~$n$.

  Comme $\C$ est algébriquement clos, $f$ possède une valeur propre
  $\lambda_1$ et un vecteur propre unitaire $\vec{v}_1$~:
  $f(\vec{v}_1) = \lambda_1 \vec{v}_1$.

  \textbf{Fait clé~:} $\vec{v}_1$ est aussi vecteur propre de $f^*$
  pour la valeur propre $\overline{\lambda_1}$. En effet,
  $\norm{(f^* - \overline{\lambda_1}\Id)(\vec{v}_1)}
  = \norm{(f - \lambda_1\Id)(\vec{v}_1)} = 0$, en utilisant le
  
ef{lem:lem:normal-norme} appliqué à l'opérateur normal
  $f - \lambda_1 \Id$.

  Posons $W = (\Vect(\vec{v}_1))^\perp$. Pour $\vec{w} \in W$~:
  $\inner{f(\vec{w}), \vec{v}_1}
  = \inner{\vec{w}, f^*(\vec{v}_1)}
  = \inner{\vec{w}, \overline{\lambda_1}\vec{v}_1}
  = \lambda_1 \inner{\vec{w}, \vec{v}_1} = 0$.
  Donc $W$ est stable par $f$.

  De même, $\inner{f^*(\vec{w}), \vec{v}_1}
  = \inner{\vec{w}, f(\vec{v}_1)}
  = \overline{\lambda_1}\inner{\vec{w}, \vec{v}_1} = 0$,
  donc $W$ est stable par $f^*$. Ainsi $g = f\restrict{W}$ est normal
  sur $W$ (de dimension $n-1$). Par hypothèse de récurrence, $W$ possède
  une base orthonormée de vecteurs propres de~$g$, et on conclut en
  ajoutant $\vec{v}_1$.
\end{proof}

% ============================================================
\section{Matrices définies positives et semi-définies positives}%
\label{sec:definie-positive}
% ============================================================

\begin{definition}{Matrice définie positive, semi-définie positive}{def:definie-pos}
  Soit $A \in \Sym_n(\R)$.\index{matrice définie positive}\index{matrice semi-définie positive}
  \begin{enumerate}[label=(\roman*)]
    \item $A$ est \emph{définie positive} si
      $\transpose{\vec{x}} A \vec{x} > 0$ pour tout
      $\vec{x} \in \R^n \setminus \set{\vzero}$.
    \item $A$ est \emph{semi-définie positive} si
      $\transpose{\vec{x}} A \vec{x} \geq 0$ pour tout $\vec{x} \in \R^n$.
    \item Les notions de \emph{définie négative} et \emph{semi-définie
      négative} s'obtiennent en inversant les inégalités.
    \item $A$ est \emph{indéfinie} si elle n'entre dans aucune des
      catégories précédentes.
  \end{enumerate}
  On note $A \succ 0$ (définie positive), $A \succeq 0$ (semi-définie
  positive).
\end{definition}

\begin{theorem}{Caractérisations des matrices définies positives}{thm:carac-def-pos}
  Pour $A \in \Sym_n(\R)$, les assertions suivantes sont
  équivalentes\index{matrice définie positive!caractérisations}~:
  \begin{enumerate}[label=(\roman*)]
    \item $A$ est définie positive~: $\transpose{\vec{x}} A \vec{x} > 0$
      pour tout $\vec{x} \neq \vzero$.
    \item Toutes les valeurs propres de $A$ sont strictement positives~:
      $\lambda_i > 0$ pour $i = 1, \ldots, n$.
    \item Tous les mineurs principaux dominants de $A$ sont strictement
      positifs\index{critère de Sylvester} (critère de Sylvester)~:
      $\Delta_k = \det(A_{1..k, 1..k}) > 0$ pour $k = 1, \ldots, n$.
    \item Il existe une matrice triangulaire inférieure $L$ à coefficients
      diagonaux strictement positifs telle que $A = L\transpose{L}$
      (décomposition de Cholesky\index{décomposition de Cholesky}).
    \item Il existe une matrice inversible $B$ telle que $A = \transpose{B}B$.
  \end{enumerate}
\end{theorem}

\begin{proof}
  Nous démontrons les implications principales.

  \medskip
  $(\mathrm{i}) \Rightarrow (\mathrm{ii})$~:
  Soit $\lambda$ une valeur propre de $A$ et $\vec{v}$ un vecteur propre
  associé ($\norm{\vec{v}} = 1$). Alors
  $0 < \transpose{\vec{v}} A \vec{v} = \lambda \transpose{\vec{v}} \vec{v}
  = \lambda$, donc $\lambda > 0$.

  \medskip
  $(\mathrm{ii}) \Rightarrow (\mathrm{i})$~:
  Par le théorème spectral, $A = P D \transpose{P}$ avec $D = \diag(\lambda_1,
  \ldots, \lambda_n)$, $\lambda_i > 0$, et $P$ orthogonale. Pour
  $\vec{x} \neq \vzero$, posons $\vec{y} = \transpose{P}\vec{x} \neq \vzero$~:
  \[
    \transpose{\vec{x}} A \vec{x}
    = \transpose{\vec{y}} D \vec{y}
    = \sum_{i=1}^{n} \lambda_i y_i^2 > 0.
  \]

  \medskip
  $(\mathrm{i}) \Rightarrow (\mathrm{iii})$~:
  Si $A \succ 0$, alors pour tout $k \leq n$, la sous-matrice
  $A_k = A_{1..k,1..k}$ est définie positive (restreindre à
  $\vec{x} = (x_1, \ldots, x_k, 0, \ldots, 0)^{\mathsf{T}}$). Donc
  ses valeurs propres sont positives, et $\Delta_k = \det(A_k) > 0$.

  \medskip
  $(\mathrm{ii}) \Rightarrow (\mathrm{iv})$~:
  Si $\lambda_i > 0$, on pose $A = P D \transpose{P}
  = P D^{1/2} \transpose{(P D^{1/2})} = \transpose{(D^{1/2}\transpose{P})}
  (D^{1/2}\transpose{P})$ avec $D^{1/2} = \diag(\sqrt{\lambda_1}, \ldots,
  \sqrt{\lambda_n})$. La factorisation de Cholesky proprement dite
  s'obtient ensuite par un algorithme constructif (élimination de Gauss
  adaptée).

  \medskip
  $(\mathrm{iv}) \Rightarrow (\mathrm{v})$~: prendre $B = L$.

  \medskip
  $(\mathrm{v}) \Rightarrow (\mathrm{i})$~:
  $\transpose{\vec{x}} A \vec{x} = \transpose{\vec{x}} \transpose{B} B \vec{x}
  = \norm{B\vec{x}}^2 \geq 0$, avec égalité si et seulement si
  $B\vec{x} = \vzero$, c'est-à-dire $\vec{x} = \vzero$ car $B$ est
  inversible.
\end{proof}

\begin{proposition}{Critère de Sylvester}{prop:sylvester}
  Soit $A \in \Sym_n(\R)$.\index{critère de Sylvester}
  \begin{enumerate}[label=(\roman*)]
    \item $A \succ 0 \iff \Delta_k > 0$ pour tout $k = 1, \ldots, n$.
    \item $A \prec 0 \iff (-1)^k \Delta_k > 0$ pour tout $k = 1, \ldots, n$.
  \end{enumerate}
  Pour la semi-définie positivité, il faut vérifier que \emph{tous} les
  mineurs principaux (pas seulement les dominants) sont $\geq 0$.
\end{proposition}

\begin{example}{Critère de Sylvester appliqué}{ex:sylvester}
  $A = \begin{pmatrix} 2 & -1 & 0 \\ -1 & 3 & -1 \\ 0 & -1 & 2 \end{pmatrix}$.

  $\Delta_1 = 2 > 0$, \quad
  $\Delta_2 = \det\begin{pmatrix} 2 & -1 \\ -1 & 3 \end{pmatrix}
  = 6 - 1 = 5 > 0$, \quad
  $\Delta_3 = \det(A) = 2(6-1) + 1(-2-0) = 10 - 2 = 8 > 0$.

  Par le critère de Sylvester, $A$ est définie positive.
\end{example}

% ============================================================
\section{Décomposition en valeurs singulières (SVD)}%
\label{sec:svd}
% ============================================================

\begin{theorem}{Décomposition en valeurs singulières}{thm:svd}
  Soit $A \in \Mat_{m \times n}(\R)$ une matrice
  quelconque.\index{décomposition en valeurs singulières}\index{SVD}
  Il existe~:
  \begin{itemize}
    \item une matrice orthogonale $U \in \Orth(m)$,
    \item une matrice orthogonale $V \in \Orth(n)$,
    \item une matrice $\Sigma \in \Mat_{m \times n}(\R)$ de la forme
      $\Sigma = \diag(\sigma_1, \ldots, \sigma_p, 0, \ldots, 0)$ avec
      $p = \min(m,n)$ et $\sigma_1 \geq \sigma_2 \geq \cdots \geq
      \sigma_p \geq 0$,
  \end{itemize}
  telles que $A = U \Sigma \transpose{V}$.

  Les $\sigma_i > 0$ sont les \emph{valeurs singulières}\index{valeurs singulières}
  de $A$~: ce sont les racines carrées des valeurs propres non nulles de
  $\transpose{A}A$ (ou de $A\transpose{A}$). Le nombre de valeurs
  singulières non nulles est égal au rang de~$A$.
\end{theorem}

\begin{proof}[Idée de la preuve]
  La matrice $\transpose{A}A \in \Sym_n(\R)$ est semi-définie positive.
  Par le théorème spectral, $\transpose{A}A = V D \transpose{V}$ avec
  $V$ orthogonale et $D = \diag(\sigma_1^2, \ldots, \sigma_r^2,
  0, \ldots, 0)$, $\sigma_i > 0$.

  Pour $i = 1, \ldots, r$, posons $\vec{u}_i = \frac{1}{\sigma_i} A \vec{v}_i$
  où les $\vec{v}_i$ sont les colonnes de $V$. On vérifie que les
  $\vec{u}_i$ sont orthonormés~:
  \[
    \inner{\vec{u}_i, \vec{u}_j}
    = \frac{1}{\sigma_i \sigma_j} \transpose{\vec{v}_i} \transpose{A} A \vec{v}_j
    = \frac{\sigma_j^2}{\sigma_i \sigma_j} \delta_{ij}
    = \delta_{ij}.
  \]

  On complète $(\vec{u}_1, \ldots, \vec{u}_r)$ en une base orthonormée
  $(\vec{u}_1, \ldots, \vec{u}_m)$ de $\R^m$ et on pose
  $U = (\vec{u}_1 \mid \cdots \mid \vec{u}_m)$. Alors $A \vec{v}_i
  = \sigma_i \vec{u}_i$ pour $i \leq r$ et $A\vec{v}_i = \vzero$ pour
  $i > r$, ce qui donne $A V = U \Sigma$, soit $A = U \Sigma \transpose{V}$.
\end{proof}

\begin{remark}{Interprétation géométrique de la SVD}{rem:svd-geo}
  L'écriture $A = U \Sigma \transpose{V}$ décompose l'application
  linéaire $\vec{x} \mapsto A\vec{x}$ en trois étapes~:
  \begin{enumerate}[label=(\arabic*)]
    \item rotation/réflexion dans $\R^n$ (via $\transpose{V}$),
    \item dilatation le long des axes et éventuellement changement de
      dimension (via $\Sigma$),
    \item rotation/réflexion dans $\R^m$ (via $U$).
  \end{enumerate}
  Géométriquement, l'image de la sphère unité de $\R^n$ par $A$ est un
  ellipsoïde dans $\R^m$ dont les demi-axes ont pour longueurs
  $\sigma_1, \ldots, \sigma_r$.
\end{remark}

\begin{example}{SVD d'une matrice $2 \times 2$}{ex:svd-2x2}
  Soit $A = \begin{pmatrix} 3 & 0 \\ 0 & 1 \end{pmatrix}$.

  $\transpose{A}A = \begin{pmatrix} 9 & 0 \\ 0 & 1 \end{pmatrix}$~:
  valeurs propres $9, 1$, donc $\sigma_1 = 3$, $\sigma_2 = 1$.
  Ici $V = I_2$, $U = I_2$, $\Sigma = A$, et la SVD est triviale.

  Pour un exemple non trivial, soit
  $B = \begin{pmatrix} 1 & 1 \\ 0 & 0 \end{pmatrix}$.

  $\transpose{B}B = \begin{pmatrix} 1 & 1 \\ 1 & 1 \end{pmatrix}$~:
  valeurs propres $2, 0$, donc $\sigma_1 = \sqrt{2}$, $\rg(B) = 1$.

  Vecteur propre de $\transpose{B}B$ pour $\lambda = 2$~:
  $\vec{v}_1 = \frac{1}{\sqrt{2}}\begin{pmatrix} 1 \\ 1 \end{pmatrix}$,
  et pour $\lambda = 0$~:
  $\vec{v}_2 = \frac{1}{\sqrt{2}}\begin{pmatrix} -1 \\ 1 \end{pmatrix}$.

  $\vec{u}_1 = \frac{1}{\sqrt{2}} B \vec{v}_1
  = \frac{1}{\sqrt{2}} \cdot \frac{1}{\sqrt{2}}\begin{pmatrix} 2 \\ 0 \end{pmatrix}
  = \begin{pmatrix} 1 \\ 0 \end{pmatrix}$.

  On complète~: $\vec{u}_2 = \begin{pmatrix} 0 \\ 1 \end{pmatrix}$.

  Donc~: $B = \begin{pmatrix} 1 & 0 \\ 0 & 1 \end{pmatrix}
  \begin{pmatrix} \sqrt{2} & 0 \\ 0 & 0 \end{pmatrix}
  \frac{1}{\sqrt{2}}\begin{pmatrix} 1 & -1 \\ 1 & 1 \end{pmatrix}^{\mathsf{T}}$.
\end{example}

% ============================================================
\section{Applications}\label{sec:spectral-applications}
% ============================================================

\subsection{Analyse en composantes principales (ACP)}%
\label{subsec:acp}

L'\emph{analyse en composantes principales}\index{ACP}\index{analyse en composantes principales}
est la technique fondamentale de réduction de dimension en statistiques.

Soit $X \in \Mat_{n \times p}(\R)$ une matrice de données centrées
($n$ observations, $p$ variables). La matrice de covariance empirique est
\[
  C = \frac{1}{n-1}\transpose{X}X \in \Sym_p(\R).
\]
Cette matrice est semi-définie positive. Par le théorème spectral,
$C = P D \transpose{P}$ avec $D = \diag(\lambda_1, \ldots, \lambda_p)$,
$\lambda_1 \geq \cdots \geq \lambda_p \geq 0$.

Les colonnes de $P$ sont les \emph{directions principales}~: la première
composante principale est la direction de variance maximale ($\lambda_1$),
la deuxième est la direction de variance maximale orthogonale à la première
($\lambda_2$), et ainsi de suite.

La proportion de variance expliquée par les $k$ premières composantes est
\[
  \frac{\sum_{i=1}^{k} \lambda_i}{\sum_{i=1}^{p} \lambda_i} = \frac{\sum_{i=1}^{k} \lambda_i}{\tr(C)}.
\]

\subsection{Classification des formes quadratiques}%
\label{subsec:formes-quad}

Une \emph{forme quadratique}\index{forme quadratique!classification}
$q \colon \R^n \to \R$ est définie par
$q(\vec{x}) = \transpose{\vec{x}} A \vec{x}$ pour une matrice symétrique
$A$ unique. Par le théorème spectral, dans la base propre orthonormée de~$A$~:
\[
  q(\vec{x}) = \lambda_1 y_1^2 + \lambda_2 y_2^2 + \cdots + \lambda_n y_n^2
\]
où $\vec{y} = \transpose{P}\vec{x}$. La classification est alors immédiate~:
\begin{itemize}
  \item \emph{Définie positive}~: tous les $\lambda_i > 0$.
  \item \emph{Semi-définie positive}~: tous les $\lambda_i \geq 0$.
  \item \emph{Définie négative}~: tous les $\lambda_i < 0$.
  \item \emph{Semi-définie négative}~: tous les $\lambda_i \leq 0$.
  \item \emph{Indéfinie}~: $\lambda_i > 0$ et $\lambda_j < 0$ pour
    certains $i, j$.
\end{itemize}
La \emph{signature}\index{signature d'une forme quadratique}
$(p, q)$ de la forme (nombre de valeurs propres positives, négatives)
est un invariant (loi d'inertie de Sylvester\index{loi d'inertie de Sylvester}).

\subsection{Optimisation~: min/max des formes quadratiques}%
\label{subsec:minmax}

\begin{theorem}{Principe du min-max (Courant--Fischer)}{thm:courant-fischer}
  Soit $A \in \Sym_n(\R)$ avec valeurs propres
  $\lambda_1 \leq \lambda_2 \leq \cdots \leq \lambda_n$.\index{Courant--Fischer}
  Alors~:
  \[
    \lambda_k = \min_{\substack{V \leq \R^n \\ \dim V = k}}
    \max_{\substack{\vec{x} \in V \\ \norm{\vec{x}} = 1}}
    \transpose{\vec{x}} A \vec{x}
    = \max_{\substack{V \leq \R^n \\ \dim V = n-k+1}}
    \min_{\substack{\vec{x} \in V \\ \norm{\vec{x}} = 1}}
    \transpose{\vec{x}} A \vec{x}.
  \]
  En particulier~:
  \[
    \lambda_1 = \min_{\norm{\vec{x}}=1} \transpose{\vec{x}} A \vec{x},
    \qquad
    \lambda_n = \max_{\norm{\vec{x}}=1} \transpose{\vec{x}} A \vec{x}.
  \]
\end{theorem}

\begin{proof}
  Démontrons le cas particulier $\lambda_n = \max_{\norm{\vec{x}}=1}
  \transpose{\vec{x}} A \vec{x}$.

  Par le théorème spectral, $A = P D \transpose{P}$ avec
  $D = \diag(\lambda_1, \ldots, \lambda_n)$. Posons $\vec{y} = \transpose{P}\vec{x}$~;
  $\norm{\vec{x}} = 1$ équivaut à $\norm{\vec{y}} = 1$, et
  \[
    \transpose{\vec{x}} A \vec{x}
    = \sum_{i=1}^{n} \lambda_i y_i^2
    \leq \lambda_n \sum_{i=1}^{n} y_i^2
    = \lambda_n.
  \]
  Le maximum est atteint pour $\vec{y} = \vece{n}$, c'est-à-dire
  $\vec{x} = P\vece{n}$ (le vecteur propre associé à $\lambda_n$).
  Le cas de $\lambda_1$ est analogue.
\end{proof}

\subsection{Vibrations et modes propres}\label{subsec:vibrations}

Considérons un système mécanique\index{vibrations!modes propres}
linéarisé au voisinage d'un équilibre~:
\[
  M \ddot{\vec{x}} + K \vec{x} = \vzero,
\]
où $M \succ 0$ (matrice de masse) et $K \succeq 0$ (matrice de raideur)
sont symétriques. On cherche des solutions harmoniques
$\vec{x}(t) = \vec{v} \cos(\omega t)$, ce qui conduit au problème aux
valeurs propres généralisé~:
\[
  K \vec{v} = \omega^2 M \vec{v}.
\]
En posant $\tilde{K} = M^{-1/2} K M^{-1/2}$ (qui est symétrique et
semi-définie positive), on se ramène au problème standard
$\tilde{K}\vec{w} = \omega^2 \vec{w}$, $\vec{w} = M^{1/2}\vec{v}$.
Par le théorème spectral, il existe une base orthonormée de
\emph{modes propres}, chacun vibrant à sa propre fréquence
$\omega_i = \sqrt{\lambda_i}$.

% ============================================================
\section{Illustrations}\label{sec:spectral-tikz}
% ============================================================

\subsection{Classification des coniques par les valeurs propres}

\begin{figure}[H]
\centering
\begin{tikzpicture}[scale=0.9]
  % Ellipse: both eigenvalues positive
  \begin{scope}[shift={(-5,0)}]
    \draw[thick,thmcolor!80,->] (-2.5,0) -- (2.5,0) node[right]{$x$};
    \draw[thick,thmcolor!80,->] (0,-2) -- (0,2) node[above]{$y$};
    \draw[thick,defcolor,rotate=30] (0,0) ellipse (1.8 and 1);
    \draw[defcolor!60,->,thick,rotate=30] (0,0) -- (1.8,0)
      node[pos=0.5,below,rotate=30]{\small $1/\sqrt{\lambda_1}$};
    \draw[defcolor!60,->,thick,rotate=30] (0,0) -- (0,1)
      node[pos=0.5,left,rotate=30]{\small $1/\sqrt{\lambda_2}$};
    \node[below] at (0,-2.3) {\small $\lambda_1, \lambda_2 > 0$ : ellipse};
  \end{scope}

  % Hyperbola: eigenvalues of opposite sign
  \begin{scope}[shift={(0,0)}]
    \draw[thick,thmcolor!80,->] (-2.5,0) -- (2.5,0) node[right]{$x$};
    \draw[thick,thmcolor!80,->] (0,-2) -- (0,2) node[above]{$y$};
    \draw[thick,defcolor,domain=-1.3:1.3,samples=100,rotate=20]
      plot ({cosh(\x)},{sinh(\x)});
    \draw[thick,defcolor,domain=-1.3:1.3,samples=100,rotate=20]
      plot ({-cosh(\x)},{sinh(\x)});
    \draw[defcolor!40,dashed,rotate=20] (-2,-2) -- (2,2);
    \draw[defcolor!40,dashed,rotate=20] (-2,2) -- (2,-2);
    \node[below] at (0,-2.3) {\small $\lambda_1 > 0 > \lambda_2$ : hyperbole};
  \end{scope}

  % Parabola-like: one eigenvalue zero
  \begin{scope}[shift={(5,0)}]
    \draw[thick,thmcolor!80,->] (-2.5,0) -- (2.5,0) node[right]{$x$};
    \draw[thick,thmcolor!80,->] (0,-2) -- (0,2) node[above]{$y$};
    \draw[thick,defcolor,domain=-1.8:1.8,samples=100]
      plot ({\x*\x/2 - 1},{\x});
    \node[below] at (0,-2.3) {\small $\lambda_1 > 0,\; \lambda_2 = 0$ : parabole};
  \end{scope}
\end{tikzpicture}
\caption{Classification des coniques
  $\transpose{\vec{x}}A\vec{x} + \vec{b}^{\mathsf{T}}\vec{x} + c = 0$
  selon les valeurs propres de la partie
  quadratique~$A$.}\label{fig:coniques}
\end{figure}

\subsection{Géométrie de la SVD}

\begin{figure}[H]
\centering
\begin{tikzpicture}[scale=1.1]
  % Left: unit circle
  \begin{scope}[shift={(-4,0)}]
    \draw[thick,thmcolor!80,->] (-1.8,0) -- (1.8,0) node[right]{$x_1$};
    \draw[thick,thmcolor!80,->] (0,-1.8) -- (0,1.8) node[above]{$x_2$};
    \draw[thick,defcolor] (0,0) circle (1);
    \draw[defcolor!70,->,very thick] (0,0) -- (1,0)
      node[below right]{\small $\vec{v}_1$};
    \draw[defcolor!70,->,very thick] (0,0) -- (0,1)
      node[above left]{\small $\vec{v}_2$};
    \node[below] at (0,-2.0) {\small Cercle unité dans $\R^2$};
  \end{scope}

  % Arrow
  \draw[-{Stealth},ultra thick,thmcolor!60] (-1.8,0) -- (-0.5,0)
    node[midway,above]{\small $A = U\Sigma \transpose{V}$};

  % Right: ellipse
  \begin{scope}[shift={(3,0)}]
    \draw[thick,thmcolor!80,->] (-2.8,0) -- (2.8,0) node[right]{$y_1$};
    \draw[thick,thmcolor!80,->] (0,-1.8) -- (0,1.8) node[above]{$y_2$};
    \draw[thick,defcolor,rotate=25] (0,0) ellipse (2.2 and 0.8);
    \draw[lemcolor,->,very thick,rotate=25] (0,0) -- (2.2,0)
      node[pos=0.6,below,rotate=25]{\small $\sigma_1 \vec{u}_1$};
    \draw[lemcolor,->,very thick,rotate=25] (0,0) -- (0,0.8)
      node[pos=0.7,left,rotate=25]{\small $\sigma_2 \vec{u}_2$};
    \node[below] at (0,-2.0) {\small Image~: ellipse de demi-axes $\sigma_1, \sigma_2$};
  \end{scope}
\end{tikzpicture}
\caption{La SVD transforme le cercle unité en un ellipsoïde. Les
  directions principales sont données par les vecteurs singuliers
  gauches $\vec{u}_i$, et les longueurs des demi-axes sont les valeurs
  singulières $\sigma_i$.}\label{fig:svd-geo}
\end{figure}

% ============================================================
\section{Exercices}\label{sec:spectral-exercices}
% ============================================================

\begin{exercise}{Diagonalisation orthogonale $2 \times 2$ \basic}{exo:diag2x2}
  Diagonaliser orthogonalement la matrice
  $A = \begin{pmatrix} 3 & 1 \\ 1 & 3 \end{pmatrix}$.
  Vérifier que $A = P D \transpose{P}$.
\end{exercise}

\begin{exercise}{Diagonalisation orthogonale $3 \times 3$ \basic}{exo:diag3x3}
  Diagonaliser orthogonalement la matrice
  $A = \begin{pmatrix} 1 & 0 & 0 \\ 0 & 0 & 1 \\ 0 & 1 & 0 \end{pmatrix}$.
\end{exercise}

\begin{exercise}{Matrice hermitienne \intermediate}{exo:hermitienne}
  Diagonaliser unitairement la matrice hermitienne
  $H = \begin{pmatrix} 2 & 1+i \\ 1-i & 3 \end{pmatrix}$.
  Vérifier que les valeurs propres sont réelles et que les vecteurs propres
  sont orthogonaux.
\end{exercise}

\begin{exercise}{Critère de Sylvester \basic}{exo:sylvester}
  En utilisant le critère de Sylvester, déterminer la nature
  (définie positive, définie négative, indéfinie) de~:
  \[
    A = \begin{pmatrix} 4 & 2 \\ 2 & 3 \end{pmatrix}, \qquad
    B = \begin{pmatrix} -3 & 1 \\ 1 & -5 \end{pmatrix}, \qquad
    C = \begin{pmatrix} 1 & 2 \\ 2 & 1 \end{pmatrix}.
  \]
\end{exercise}

\begin{exercise}{Formes quadratiques et coniques \intermediate}{exo:conique}
  Classifier la conique $5x^2 + 4xy + 2y^2 = 6$ en déterminant les
  valeurs propres de la matrice associée à la forme quadratique.
  Tracer la conique dans la base propre.
\end{exercise}

\begin{exercise}{Calcul de SVD \intermediate}{exo:svd-calcul}
  Calculer la décomposition en valeurs singulières de la matrice
  $A = \begin{pmatrix} 1 & 1 \\ 1 & -1 \\ 0 & 0 \end{pmatrix}$.
\end{exercise}

\begin{exercise}{Décomposition de Cholesky \intermediate}{exo:cholesky}
  Calculer la décomposition de Cholesky de la matrice
  $A = \begin{pmatrix} 4 & 2 & -2 \\ 2 & 5 & -4 \\ -2 & -4 & 14 \end{pmatrix}$.
  Vérifier que $A = L\transpose{L}$.
\end{exercise}

\begin{exercise}{Puissances d'une matrice symétrique \intermediate}{exo:puissance-sym}
  Soit $A = \begin{pmatrix} 5 & 4 \\ 4 & 5 \end{pmatrix}$.
  \begin{enumerate}[label=(\alph*)]
    \item Diagonaliser orthogonalement~$A$.
    \item Calculer $A^{10}$.
    \item Exprimer $e^{tA} = \sum_{k=0}^{\infty} \frac{t^k}{k!} A^k$
      sous forme explicite.
  \end{enumerate}
\end{exercise}

\begin{exercise}{Opérateur normal non auto-adjoint \intermediate}{exo:normal}
  Soit $A = \begin{pmatrix} 1 & -1 \\ 1 & 1 \end{pmatrix} \in \Mat_2(\R)$.
  \begin{enumerate}[label=(\alph*)]
    \item Montrer que $A$ n'est pas symétrique.
    \item Vérifier que $A$ est normale~: $\transpose{A}A = A\transpose{A}$.
    \item Montrer que $A$ n'est pas diagonalisable sur $\R$ mais
      l'est sur $\C$. Trouver une matrice unitaire $U$ telle que
      $A = U D \hermit{U}$.
  \end{enumerate}
\end{exercise}

\begin{exercise}{Approximation de rang faible et SVD \challenging}{exo:approx-rang}
  Soit $A \in \Mat_{m \times n}(\R)$ de SVD $A = U \Sigma \transpose{V}$.
  \begin{enumerate}[label=(\alph*)]
    \item Montrer que pour tout $k \leq \rg(A)$, la matrice
      $A_k = \sum_{i=1}^{k} \sigma_i \vec{u}_i \transpose{\vec{v}_i}$
      est de rang~$k$.
    \item Admettre le théorème d'Eckart--Young\index{théorème d'Eckart--Young}~:
      $A_k$ est la meilleure approximation de rang~$k$ de~$A$ au sens de
      la norme de Frobenius~:
      $\norm{A - A_k}_F = \min_{\rg(B) \leq k} \norm{A - B}_F$.
    \item Calculer $\norm{A - A_k}_F^2$ en fonction des valeurs singulières.
    \item Application~: si $A = \begin{pmatrix} 3 & 0 \\ 0 & 2 \\ 0 & 0
      \end{pmatrix}$, quelle est la meilleure approximation de rang~$1$~?
  \end{enumerate}
\end{exercise}

\begin{exercise}{Rayleigh et convergence \challenging}{exo:rayleigh}
  Soit $A \in \Sym_n(\R)$ avec valeurs propres
  $\lambda_1 < \lambda_2 \leq \cdots \leq \lambda_n$. Le
  \emph{quotient de Rayleigh}\index{quotient de Rayleigh} est
  $R(\vec{x}) = \frac{\transpose{\vec{x}} A \vec{x}}{\transpose{\vec{x}} \vec{x}}$
  pour $\vec{x} \neq \vzero$.
  \begin{enumerate}[label=(\alph*)]
    \item Montrer que $\lambda_1 \leq R(\vec{x}) \leq \lambda_n$.
    \item Montrer que $R(\vec{x}) = \lambda_k$ si et seulement si $\vec{x}$
      est vecteur propre de $A$ pour $\lambda_k$.
    \item Calculer le gradient $\nabla R(\vec{x})$ et montrer que ses
      points critiques sont les vecteurs propres de~$A$.
  \end{enumerate}
\end{exercise}

\begin{exercise}{Preuve de la décomposition de Cholesky \challenging}{exo:cholesky-preuve}
  Soit $A \in \Sym_n(\R)$ définie positive.
  \begin{enumerate}[label=(\alph*)]
    \item Montrer que $a_{11} > 0$.
    \item Écrire $A$ sous la forme par blocs
      $A = \begin{pmatrix} a_{11} & \transpose{\vec{c}} \\
      \vec{c} & A' \end{pmatrix}$ et montrer que
      $S = A' - \frac{1}{a_{11}} \vec{c}\transpose{\vec{c}}$ (le
      complément de Schur\index{complément de Schur}) est définie
      positive de taille $(n-1) \times (n-1)$.
    \item En déduire par récurrence l'existence de la décomposition
      $A = L\transpose{L}$ avec $L$ triangulaire inférieure à diagonale
      strictement positive.
    \item Montrer l'unicité de cette décomposition.
  \end{enumerate}
\end{exercise}

% ============================================================
\section*{Résumé du chapitre}
\addcontentsline{toc}{section}{Résumé du chapitre}
% ============================================================

\begin{tcolorbox}[colback=defcolor!3,colframe=defcolor!60,
  title={\textbf{Résumé --- Chapitre \thechapter}},
  fonttitle=\large,breakable]

\textbf{Opérateurs auto-adjoints.}
$f$ est auto-adjoint si $f = f^*$, i.e.\
$\inner{f(\vec{u}),\vec{v}} = \inner{\vec{u},f(\vec{v})}$ pour tous
$\vec{u}, \vec{v}$. En base orthonormée~: matrice symétrique ($\R$) ou
hermitienne ($\C$).

\medskip
\textbf{Propriétés spectrales fondamentales.}
\begin{itemize}[nosep]
  \item Les valeurs propres d'un opérateur auto-adjoint sont \emph{réelles}.
  \item Les sous-espaces propres associés à des valeurs propres distinctes
    sont \emph{orthogonaux}.
\end{itemize}

\medskip
\textbf{Théorème spectral (réel).}
Toute matrice symétrique réelle $A \in \Sym_n(\R)$ est
\emph{orthogonalement diagonalisable}~:
$A = P D \transpose{P}$, $P \in \Orth(n)$, $D$ diagonale réelle.

\medskip
\textbf{Théorème spectral (hermitien).}
Toute matrice hermitienne $A = \hermit{A}$ est
\emph{unitairement diagonalisable}~:
$A = U D \hermit{U}$, $U$ unitaire, $D$ diagonale réelle.

\medskip
\textbf{Opérateurs normaux.}
$f$ est normal si $ff^* = f^*f$. Sur $\C$, $f$ est normal $\iff$
$f$ admet une base orthonormée de vecteurs propres.

\medskip
\textbf{Matrices définies positives.}
$A \succ 0$ $\iff$ $\lambda_i > 0\ \forall i$ $\iff$ $\Delta_k > 0\
\forall k$ (Sylvester) $\iff$ $A = L\transpose{L}$ (Cholesky) $\iff$
$A = \transpose{B}B$, $B$ inversible.

\medskip
\textbf{SVD.}
Toute matrice $A \in \Mat_{m \times n}(\R)$ s'écrit
$A = U \Sigma \transpose{V}$ avec $U \in \Orth(m)$, $V \in \Orth(n)$,
$\Sigma = \diag(\sigma_1, \ldots, \sigma_p)$, $\sigma_i \geq 0$.
Les $\sigma_i$ sont les racines carrées des valeurs propres de
$\transpose{A}A$.

\medskip
\textbf{Applications.}
ACP (directions de variance maximale), classification des formes
quadratiques et coniques, principe du min-max (Courant--Fischer),
approximation de rang faible (Eckart--Young), vibrations et modes propres.
\end{tcolorbox}
